\section{Experimental Verifications: Testing the TRXT Phase Transition}
\label{sec:predictions}

The physical core of the TRXT model is the identification of the vacuum as a Superfluid Condensate. This implies a cosmological Phase Transition effectively occurring at the condensation scale.
If this transition corresponds to the Recombination Era ($z_c \approx 1100, T_c \approx 1 \text{ eV}$), as hypothesized to resolve the Coincidence Problem, it leads to specific, falsifiable signatures.

\subsection{CMB Acoustic Scale Shifts}
We modeled the impact of a TRXT phase transition injecting a vacuum energy component $\Delta \Omega_{vac}$ around $z=1100$ (Early Dark Energy scenario).
The presence of this field increases the Hubble rate $H(z)$ prior to recombination, thereby shrinking the sound horizon $r_s$:
\begin{equation}
    r_s = \int_{z_{rec}}^\infty \frac{c_s dz}{H(z)}
\end{equation}
Our simulations (`v16_cmb_predictions.py`) show that an energy injection of 10\% of the total density results in a $-0.65\%$ shift in the sound horizon.
To maintain the observed angular size $\theta_* = r_s / D_A$, the Hubble Constant $H_0$ would need to increase by $\approx 0.6\%$.
\textbf{Conclusion:} A Recombination-era transition is consistent with Planck data only if the energy fraction is small ($< 5\%$). It does not, by itself, fully resolve the Hubble Tension (which requires a $\sim 9\%$ shift) unless the transition has broader temporal width.

\textit{Important Note on Scales:} This Recombination-era transition ($T_c \sim 1 \text{ eV}$) is distinct from the \textbf{Primordial Condensation} ($T_c \sim 10^{15} \text{ GeV}$). It represents a secondary vacuum realignment, physically analogous to the \textbf{Helium-3 A-to-B Phase Transition} in condensed matter systems. While the Primordial transition established spacetime itself, this secondary transition corresponds to the \textbf{Early Dark Energy (EDE)} mechanism proposed to resolve the Hubble Tension ($5\sigma$ discrepancy). Thus, far from being an ad-hoc addition, this "Late Transition" is a requisite feature of superfluid vacuum models and provides a natural explanation for current cosmological anomalies.



\subsection{Gravitational Wave Background}
A first-order cosmological phase transition generates a stochastic background of gravitational waves from bubble collisions.
For a transition at $T_* = 1 \text{ eV}$, the peak frequency of these waves today is redshifted to:
\begin{equation}
    f_{peak} \approx 10^{-14} \text{ Hz}
\end{equation}
This extremely low frequency lies far below the sensitivity of LISA ($10^{-3}$ Hz) or SKA ($10^{-9}$ Hz). However, these waves have wavelengths comparable to the Hubble radius at recombination, meaning they would generate Tensor Modes in the CMB polarization.
Our calculation (`v16_gw_predictions.py`) for a transition strength $\alpha = 0.05$ predicts a peak amplitude:
\begin{equation}
    \Omega_{GW} h^2 \approx 7 \times 10^{-13}
\end{equation}
This exceeds the current upper limits on the tensor-to-scalar ratio $r$ from Planck/BICEP ($\Omega_{GW} \lesssim 10^{-16}$).
\textbf{Constraint:} This result \textbf{rules out} a strong first-order transition ($\alpha > 0.005$) at the eV scale. The TRXT transition must be either a **Second-Order (smooth) crossover** or sufficiently weak to satisfy B-mode constraints. 
\textit{Clarification for V7.0:} "Early Phase Transition" refers to the timing ($z \sim 1100$), not the violence. The data requires an early but gentle restructuring of the vacuum, akin to a second-order magnetic transition rather than a first-order boiling process.
This provides a precise experimental bound on the theory's parameter space.

